\chapter{Mathematical contributions and relation to existing theory}\label{app:contrib}

This appendix places the results of the volume in their mathematical setting. The techniques are classical and the objects are new: the contribution of the UNNS theory is the identification of the collective map \(C\), of the coordinates (\(\gamma\), exponent state, absorption sets) in which its dynamics become finite, and the theorems that follow. The table pairs each group of results with the established background it uses.

\footnotesize
\begin{longtable}{@{}p{34mm}p{36mm}p{52mm}@{}}
\toprule Results & Established tools used & What the theory establishes\\\midrule\endfirsthead
\toprule Results & Established tools used & What the theory establishes\\\midrule\endhead
Fibre growth (Thms~\ref{thm:growth-time},~\ref{thm:growth}, Prop.~\ref{prop:rate}) & valuations, gcd/lcm, divisor counting & unbounded multiplicity of the collective orbit, an explicit growth-time bound, and a cubic-root lower rate\\
Regime~I (Thms~\ref{thm:regimeI},~\ref{thm:L3},~\ref{thm:L4}) & squarefree decomposition; iteration of \(t\mapsto\mathrm{sqf}(t+1)\) & closed dynamics at \(G=1\), exact three-step scaling and exact reach counting\\
Exclusions (Thms~\ref{thm:G236},~\ref{thm:G24}, Lem.~\ref{lem:escape}) & residues; descent on positive integers & \(G_\infty\) can never divide 24, by a descent argument on \(\gamma\)\\
Persistent regimes (Thms~\ref{thm:Pg},~\ref{thm:T42},~\ref{thm:Pstar}, Prop.~\ref{prop:T38a}) & orders of 2 modulo \(g\); Dirichlet's theorem (only for Prop.~\ref{prop:T38a}) & explicit infinite families with \(G_\infty=g>1\), including \(3\mid g\): universal collapse is false\\
\(\gamma\)-framework, window (Thms~\ref{thm:gamma},~\ref{thm:T47}) & integrality; congruences modulo 8, 3, 24 & an integer coordinate in which persistent orbits live in a bounded window\\
Class closure, exact scaling (Thm~\ref{thm:closure}) & finite deterministic systems are eventually periodic; structure of \((\Z/g)^\times\) & every orbit eventually satisfies \(x_{n+P}=\rho^2x_n\); the finite-state data are identified\\
Carry calculus (Thms~\ref{thm:absorb}--\ref{thm:odometer}) & cyclic subgroups, orders of primes modulo \(g\) & carry completeness, prime-power carries, the odometer law\\
\(q\)-carry family (Thm~\ref{thm:qcarry}) & cyclic quotient \(\langle q\rangle\), \(\langle2\rangle H\) & an explicit orbit family with cycle \(1,\dots,q^2-1\); realised for \(q=3,5,7\)\\
Single-prime dominance (Thms~\ref{thm:scheduler},~\ref{thm:dominance},~\ref{thm:rho}) & extremal (worst-case) argument on cyclic odometers & exactly one prime carries; \(\rho\) is a prime power; logarithmic reach (Cor.~\ref{cor:logreach})\\
Selection and certificates (Thms~\ref{thm:predictor}--\ref{thm:consequences}) & quotient groups \((\Z/g)^\times/H\); minimax & survivor decided by a finite residue test at one reset state\\
Effective bounds (Thms~\ref{thm:kappa}--\ref{thm:dynG}) & iteration of \(a\mapsto\lfloor a/4\rfloor+1\); pigeonhole on classes; epoch decomposition & computable bounds on the transient from \((g_0,\gamma_0)\)\\
\bottomrule
\end{longtable}
\normalsize

\paragraph{Established background.} The following facts are used as standard: the unit group \((\Z/p)^\times\) of a prime is cyclic, and the order and subgroup structure of a cyclic group are those of \(\Z/n\) \cite{ireland,hardy}; Dirichlet's theorem on primes in arithmetic progressions \cite{apostol}, used once (Proposition~\ref{prop:T38a}) and avoided in Theorem~\ref{thm:T42}; the pigeonhole principle for a deterministic map on a finite set; valuations and the characterisation of gcd and lcm through them. None of these is specific to the UNNS Substrate.

\paragraph{Results about UNNS structures.} What is specific to the Substrate is the following chain, each link a theorem of this volume. The collective rule \(C\) is chosen (Definition~\ref{th:F.CDEF}); once chosen, its orbits have unbounded multiplicity (Theorem~\ref{thm:growth}); the coordinate \(\gamma\) and the one-step law reduce the dynamics to residues (Theorems~\ref{thm:gamma},~\ref{thm:onestep}); every orbit is eventually exactly self-similar (Theorem~\ref{thm:closure}); the self-similar regime is a single odometer (Theorem~\ref{thm:dominance}); which odometer is decided by finite arithmetic (Theorems~\ref{thm:minimax},~\ref{thm:certificate}); and the time to reach it is bounded effectively (Theorem~\ref{thm:dynG}). The distinction between existence (eventual exact scaling), structure (one prime, one carry word) and effectiveness (computable transient bounds) is the spine of the volume.

\paragraph{Open directions.} The scope of certification beyond prime \(g\), sharper transient bounds, the classification of the \(g\) that admit persistent orbits, and the dynamics of collective rules other than \(C\) are discussed in Appendix~\ref{app:open}. The collective rule is a selected rule: Volume~I shows that the interaction axioms considered there do not select a unique rule, so every statement of this volume is a statement about \(C\).
